% BEGIN EMBEDDED BIBLIOGRAPHY
% Bibliography data are part of this source; the compiler writes its temporary input.
\begin{filecontents*}[overwrite]{\jobname.bib}
@unpublished{QED3Manuscript,
  title = {The field-strength {Osterwalder--Schrader} axioms for {QED} in three dimensions},
  date = {2026},
  shorthand = {Cur},
  note = {Comprehensive edition, 263 pp., 20 October 2026, before the modular integration revision}
}
@book{Durrett2019,
  author = {Durrett, Rick},
  title = {Probability: Theory and Examples},
  edition = {5},
  publisher = {Cambridge University Press},
  date = {2019},
  url = {https://sites.math.duke.edu/~rtd/PTE/PTE5_011119.pdf},
  note = {Author's version of 11 January 2019}
}
\end{filecontents*}
% END EMBEDDED BIBLIOGRAPHY
% ECM1. Root-owned extraction of the continuum-law theorem.
\documentclass[11pt,letterpaper,reqno]{amsart}
% BEGIN INLINE manuscript-style
\makeatletter
\NeedsTeXFormat{LaTeX2e}[2020/10/01]
% Use with \documentclass[11pt,letterpaper,reqno]{amsart}.
% This file controls presentation only. It supplies no mathematical content.

\RequirePackage[T1]{fontenc}
\RequirePackage{lmodern}
\RequirePackage{microtype}
\RequirePackage[letterpaper,margin=1in,headheight=15pt]{geometry}
\RequirePackage{amsmath,amsthm,amssymb,mathtools}
\RequirePackage{enumitem}
\RequirePackage{perpage}
\RequirePackage{xcolor}
\definecolor{manuscriptred}{rgb}{0.6,0.1,0.1}
\definecolor{manuscriptblue}{rgb}{0.1,0.3,0.5}
\definecolor{manuscriptmagenta}{HTML}{60445C}
\RequirePackage[breaklinks=true,colorlinks=true,
  linkcolor=manuscriptred,citecolor=manuscriptblue,urlcolor=manuscriptmagenta]{hyperref}
\RequirePackage{cleveref}
\hypersetup{pdfauthor={}}

% Retain the AMS paragraph, title, heading and running-head conventions.
% Equations retain one global Arabic sequence, as in the reference.
\setcounter{tocdepth}{1}
\renewcommand{\thefigure}{\Roman{figure}}
\renewcommand{\thetable}{\Roman{table}}
\MakePerPage{footnote}
\renewcommand{\thefootnote}{%
  \ifnum\value{footnote}>9\arabic{footnote}\else\fnsymbol{footnote}\fi}

% Upright statements. Main results are A, B, C; supporting statements
% share a section-local sequence. These end marks indicate typography,
% never proof completeness, validation, or a deferred-proof status.
\theoremstyle{definition}
\newtheorem{ms@theorem}{Theorem}
\renewcommand{\thems@theorem}{\Alph{ms@theorem}}
\newtheorem*{ms@theoremstar}{Theorem}
\newtheorem{ms@proposition}{Proposition}[section]
\newtheorem{ms@lemma}[ms@proposition]{Lemma}
\newtheorem{ms@definition}[ms@proposition]{Definition}
\newtheorem{ms@corollary}[ms@proposition]{Corollary}
\newtheorem*{ms@corollarystar}{Corollary}

\newcommand{\ms@majorstart}{%
  \pushQED{\qed}\renewcommand{\qedsymbol}{$\blacksquare\square$}}
\newcommand{\ms@minorstart}{%
  \pushQED{\qed}\renewcommand{\qedsymbol}{$\blacksquare$}}
\newenvironment{theorem}[1][]{\ms@majorstart\ms@theorem[#1]}
  {\popQED\endms@theorem}
\newenvironment{theorem*}[1][]{\ms@majorstart\ms@theoremstar[#1]}
  {\popQED\endms@theoremstar}
\newenvironment{proposition}[1][]{\ms@minorstart\ms@proposition[#1]}
  {\popQED\endms@proposition}
\newenvironment{lemma}[1][]{\ms@minorstart\ms@lemma[#1]}
  {\popQED\endms@lemma}
\newenvironment{definition}[1][]{\ms@minorstart\ms@definition[#1]}
  {\popQED\endms@definition}
\newenvironment{corollary}[1][]{\ms@majorstart\ms@corollary[#1]}
  {\popQED\endms@corollary}
\newenvironment{corollary*}[1][]{\ms@majorstart\ms@corollarystar[#1]}
  {\popQED\endms@corollarystar}
% The source's lemma/proposition variants are deliberately not carried over:
% choose one marker for each statement type rather than encode proof status.

\theoremstyle{remark}
\newtheorem{remark}[ms@proposition]{Remark}
\newtheorem*{remark*}{Remark}
% The ordinary amsthm proof environment retains its open square.

\crefname{ms@theorem}{theorem}{theorems}
\Crefname{ms@theorem}{Theorem}{Theorems}
\crefname{ms@proposition}{proposition}{propositions}
\Crefname{ms@proposition}{Proposition}{Propositions}
\crefname{ms@lemma}{lemma}{lemmas}
\Crefname{ms@lemma}{Lemma}{Lemmas}
\crefname{ms@definition}{definition}{definitions}
\Crefname{ms@definition}{Definition}{Definitions}
\crefname{ms@corollary}{corollary}{corollaries}
\Crefname{ms@corollary}{Corollary}{Corollaries}
\crefname{remark}{remark}{remarks}
\Crefname{remark}{Remark}{Remarks}

% Bibliography configuration for this manuscript.
  \RequirePackage[backend=biber,style=alphabetic,backref=true,
    giveninits=true,isbn=false,date=year]{biblatex}
  \addbibresource{\jobname.bib}
  \setcounter{biburlnumpenalty}{100}
  \setcounter{biburllcpenalty}{100}
  \setcounter{biburlucpenalty}{100}
  \renewcommand*{\bibfont}{\footnotesize}
  \renewbibmacro{in:}{}
  \DeclareFieldFormat{pages}{#1}
  % Use the public alphabetic style. No private author-hash overrides;
  % repeated author names are printed in full.
\makeatother
% END INLINE manuscript-style
\newcommand{\R}{\mathbb R}
\newcommand{\C}{\mathbb C}
\newcommand{\Q}{\mathbb Q}
\newcommand{\E}{\mathbb E}
\newcommand{\dd}{\,\mathrm d}
\newcommand{\Sspace}{\mathcal S}
\title[Ergodic components and clustering]
 {Ergodic components and clustering for reflection-positive fields}
\date{}
\hypersetup{pdftitle={Ergodic components and clustering for reflection-positive fields},pdfauthor={}}
\begin{document}
\begin{abstract}
A translation-invariant law of finite-component tempered fields in
three dimensions is assumed invariant and reflection positive at the
coordinate and diagonal planes. We prove clustering relative to its
translation-invariant sigma-algebra. Almost every translation-ergodic
component retains the reflection properties and is mixing in every
direction. The proof uses a diagonal reflection to obtain a spectral
cone, then identifies zero-energy vectors with translation-invariant
observables. If the parent law also satisfies an \(L^2\) Gaussian
bound, the same constant holds in almost every component.
\end{abstract}
\maketitle
\tableofcontents
\section{Introduction}\label{sec:introduction}
Let \(V\) be a finite-dimensional real inner-product space with an
orthogonal representation \(\rho:O(3)\to O(V)\). Let
\(\mathcal X=\mathcal S'(\R^3;V)\), equipped with the sigma-algebra
generated by evaluations. For \(g(x)=Qx+b\), use
\begin{equation}\label{eq:euclidean-action}
 (g\cdot f)(x)=\rho(Q)f(Q^{-1}(x-b)),\qquad
 (g\cdot F)(f)=F(g^{-1}\cdot f),\qquad
 \widehat g A(F)=A(g^{-1}\cdot F).
\end{equation}
Write \(U(a)=\widehat{\tau_a}\).
For open \(O\subset\R^3\), let \(\Sigma_O\) be generated by
\(F(f)\) with real \(f\in C_c^\infty(O;V)\).
For a reflection \(r\), put \(\mathcal R_r=\widehat r\).
An \(r\)-invariant law \(\lambda\) is reflection positive on an
open side \(H^+\) if
\[
 \E_\lambda[\overline{\mathcal R_r A}\,A]\ge0
\]
for finite character sums \(A=\sum_bz_be^{iF(f_b)}\) with
real \(f_b\in C_c^\infty(H^+;V)\).
The positive sides below are chosen by making the displayed linear
coordinate function positive.

\begin{theorem}\label{thm:components-clustering}
Let \(\lambda\) be a translation-invariant probability law on
\(\mathcal X\), invariant and reflection positive at the nine planes
\begin{equation}\label{eq:nine-planes}
 x_j=0\quad(1\le j\le3),\qquad
 x_j-x_k=0,\quad x_j+x_k=0\quad(1\le j<k\le3).
\end{equation}
Let \(\mathcal I\) be the translation-invariant sigma-algebra,
with invariance understood modulo \(\lambda\)-null sets, and let
\(P=\E_\lambda[\,\cdot\mid\mathcal I]\). Then
\begin{equation}\label{eq:main-relative-clustering}
 \langle A,U(a)B\rangle_{L^2(\lambda)}
 \longrightarrow\langle A,PB\rangle_{L^2(\lambda)}
 \qquad(|a|\longrightarrow\infty)
\end{equation}
for every \(A,B\in L^2(\lambda)\).

There is a unique translation-ergodic decomposition
\(\lambda=\int\lambda_\xi\,\dd\nu(\xi)\), understood as the
distribution of its conditional law on \(\mathcal X\).
On one common \(\nu\)-full set, every \(\lambda_\xi\) is
translation ergodic, invariant and reflection positive at every
real translate of the planes in \eqref{eq:nine-planes}, and
\begin{equation}\label{eq:main-component-clustering}
 \langle A,U(a)B\rangle_{L^2(\lambda_\xi)}
 \longrightarrow
 \overline{\E_{\lambda_\xi}A}\,\E_{\lambda_\xi}B
 \qquad(|a|\longrightarrow\infty)
\end{equation}
for every \(A,B\in L^2(\lambda_\xi)\).
Reflection positivity holds on the entire positive-side
\(L^2(\lambda_\xi)\) space. These component laws are cubically invariant.

If in addition, for some fixed \(C\ge0\),
\begin{equation}\label{eq:probability-G}
 \E_\lambda e^{F(f)}\le e^{C\|f\|_2^2/2}
 \qquad(f\in\mathcal S(\R^3;V)\text{ real}),
\end{equation}
then the common full-measure set above may be chosen so that every
component in it satisfies \eqref{eq:probability-G} with the same
\(C\), simultaneously for all real tests.
Any specified measurable event of full \(\lambda\)-measure has
full measure in almost every component.
\end{theorem}

The proof follows the reflection and component arguments of
\cite[Section~7]{QED3Manuscript}. They depend on the law and its
geometric action, so the proof applies to the finite component
space \(V\) above. No field equation or Gaussian estimate is used
for \eqref{eq:main-relative-clustering} or the reflection properties
of the components.

The spectral step is essential for the all-direction limit.
Coordinate reflection positivity constructs a nonnegative generator
\(H_j\) and commuting transverse momenta \(\mathsf P_i\).
A second reconstruction at a diagonal mirror gives a bound on
Laplace transforms of their joint spectral measure. We use
transforms whose suprema on the positive real and imaginary axes
have bounded product, while their modulus grows exponentially
outside the sector \(E\ge|p_i|\).
This proves \(H_j\ge|\mathsf P_i|\).
At zero energy all transverse momenta therefore vanish.
It follows that invariance under one coordinate translation group
is already invariance under all translations. Semigroup convergence
then gives \eqref{eq:main-relative-clustering}, uniformly in
transverse displacements.

For the component assertions, a bounded translation-invariant
observable has a representative in either half-space and is fixed
by its reflection. Applying reflection positivity to \(\sqrt h A\),
with \(h\) a nonnegative invariant weight, yields a conditional
reflection inequality. Countable determining character algebras
put the resulting properties on one set of component laws.
The optional Gaussian assertion uses a separate estimate:
spatial averages of \(e^{F(f)}\) have moments controlled by
translated \(L^2\) overlaps of \(f\). Those overlaps vanish at
large separation, which preserves the original constant \(C\)
after conditioning.

The conclusion is qualitative clustering. It asserts no positive
spectral gap. The hypotheses also do not assert invariance under
arbitrary rotations or field-sign change in each component.
\section{Reflection quotients}\label{sec:quotients}
For an open set $O$, let $\Sigma_O$ be generated by $F(f)$ with
real $f\in C_c^\infty(O;V)$. Let $\lambda$ be a
translation-invariant law and let $r_j$ reflect in $x_j=0$.
Suppose that $\lambda$ is invariant and reflection positive at this
plane. For a Euclidean motion $g$ write
$\widehat g A(F)=A(g^{-1}F)$ and $U(a)=\widehat{\tau_a}$.
For a reflection $r$ put
$\mathcal R_r=\widehat r$. Write $\mathcal R_j=\mathcal R_{r_j}$ and
$\Theta_j A=\overline{\mathcal R_j A}$.
On $\mathcal E_j^+=L^2(\Sigma_{\{x_j>0\}},\lambda)$ the form
\[
 b_j(A,B)=\E_\lambda[(\Theta_j A)B]
\]
is positive semidefinite and $b_j(A,A)\le\|A\|_{L^2(\lambda)}^2$.
Let $\mathcal H_j$ be its quotient completion, $[A]_j$ the quotient
map, and $\Omega_j=[1]_j$.
Reflection positivity stated on the characters $e^{iF(f)}$, with
real $f\in C_c^\infty(\{x_j>0\};V)$, extends to this form:
their linear span is dense in the positive-side $L^2$ space, and the
form is bounded by the $L^2$ norm. The character-density argument is
given in \cref{subsec:component-hypotheses}.

\begin{proposition}\label{prop:axis-reconstruction}
There are a non-negative self-adjoint $H_j$ and a strongly continuous
unitary representation $V_j$ of $e_j^\perp$ on $\mathcal H_j$ such
that
\begin{equation}\label{eq:axis-semigroup}
 e^{-tH_j}V_j(y)[A]_j=[U(te_j+y)A]_j,
 \qquad t\ge0,\quad y\perp e_j.
\end{equation}
The two representations commute, and $H_j\Omega_j=0$.
\end{proposition}

\begin{proof}
Translations are strongly continuous on $L^2(\lambda)$. On
characters this follows from continuity of translation of a Schwartz
test and bounded convergence; their density gives it on all of
$L^2(\lambda)$. Positive time translations preserve the positive-side
algebra. With $|A|_b=b_j(A,A)^{1/2}$, reflection and translation
invariance give
\[
 |U(te_j)A|_b^2=b_j(A,U(2te_j)A)
 \le |A|_b\,|U(2te_j)A|_b.
\]
Iteration yields
\[
 |U(te_j)A|_b
 \le |A|_b^{1-2^{-k}}\|A\|_{L^2(\lambda)}^{2^{-k}}.
\]
Letting $k\to\infty$ proves that $[A]_j\mapsto[U(te_j)A]_j$
is a well-defined contraction. It is symmetric, obeys the semigroup
law, and is strongly continuous by the contraction of the quotient
map. Writing it as $T_j(t)$, one has
$T_j(t)=T_j(t/2)^2\ge0$. If $T_j(t)\psi=0$, repeated square roots
and strong continuity give $\psi=0$. Spectral calculus therefore
defines $H_j=-\log T_j(1)\ge0$; uniqueness of positive roots gives
$T_j(t)=e^{-tH_j}$ first at dyadic times and then at all times.

For $y\perp e_j$, $U(y)$ preserves the half-space and commutes with
reflection, so it is an isometry for $b_j$, with inverse $U(-y)$.
Its quotient extension is the unitary $V_j(y)$. Strong continuity
and commutation with $T_j(t)$ follow from the corresponding Euclidean
translation identities. The constant observable is fixed throughout.
\end{proof}

Write $V_j(se_i)=e^{is\mathsf P_i}$ for $i\ne j$. The joint spectral
measure $\mathsf E_j$ of $H_j$ and these commuting momenta lies on
$[0,\infty)\times e_j^\perp$. For $\psi\in\mathcal H_j$ let
$\varrho_\psi(A)=\langle\psi,\mathsf E_j(A)\psi\rangle$.
Then
\begin{equation}\label{eq:joint-spectral-integral}
 \langle\psi,e^{-tH_j}V_j(y)\psi\rangle
 =\int e^{-tE+ip\cdot y}\,\dd\varrho_\psi(E,p).
\end{equation}

\section{The spectral cone from a diagonal reflection}\label{sec:cone}

\begin{proposition}\label{prop:diagonal-cone}
Let $i\ne j$, and suppose in addition that $\lambda$ is invariant
and reflection positive at $x_j+x_i=0$. Then
\begin{equation}\label{eq:sharp-coordinate-cone}
 \mathsf E_j\{(E,p):E<|p_i|\}=0.
\end{equation}
Thus $H_j\ge|\mathsf P_i|$ in the joint spectral calculus.
\end{proposition}

\begin{proof}
Fix $D>0$ and $A\in L^2(\Sigma_B,\lambda)$, where
$B=\{0<x_j<D,\ |x_i|<D\}$, and set $\psi=[A]_j$.
Let $\varrho$ be the image of $\varrho_\psi$ under
$(E,p)\mapsto(E,p_i)$. This is a finite positive measure on the
right half-plane, identified by $z=E-ip_i$.
For $a,c\ge0$ consider
\[
 K(a,c)=\E_\lambda\bigl[(\Theta_j A)
 U(2De_j+a(e_j+e_i)+c(e_j-e_i))A\bigr].
\]
The axis reconstruction gives
\begin{equation}\label{eq:cone-first-picture}
 K(a,c)=\int e^{-2DE}e^{-az-c\overline z}\,\dd\varrho(E,p_i).
\end{equation}

There is a second representation with the diagonal as time direction.
Put $n_\pm=(e_j\pm e_i)/\sqrt2$, $y=D(e_j+e_i)/2$, and let
$r_+$ reflect in $x_j+x_i=0$. The observables
\[
 A_1=\mathcal R_{r_+}U(-y)\mathcal R_j A,
 \qquad A_2=U(2De_j-y)A
\]
are supported on the positive side of this diagonal. Indeed,
$r_jB-y$ lies in its negative side and $B+2De_j-y$ in its positive
side. Their $L^2$ norms equal $\|A\|_2$.
Writing $\Theta_+A=\overline{\mathcal R_{r_+}A}$, translation of
the original expectation by $-y$ gives
\[
 K(a,c)=\E_\lambda\bigl[(\Theta_+A_1)
 U(a(e_j+e_i)+c(e_j-e_i))A_2\bigr].
\]
Choose $Q\in O(3)$ carrying $e_j,e_i$ to $n_+,n_-$ and apply
\cref{prop:axis-reconstruction} to the transformed law
$(Q^{-1})_*\lambda$. If $\mathsf E'$ is its joint spectral measure,
set $\psi_a=[\widehat{Q^{-1}}A_a]'$, $a=1,2$, and let
$\varrho'$ be the image of
$\langle\psi_1,\mathsf E'(\cdot)\psi_2\rangle$
under $(E,p)\mapsto(\sqrt2E,\sqrt2p_i)$. This gives
\begin{equation}\label{eq:cone-second-picture}
 K(a,c)=\int e^{-aE'+icp'}\,\dd\varrho'(E',p'),
 \qquad \|\varrho'\|_{\mathrm{TV}}\le\|A\|_2^2,
\end{equation}
where $E'\ge0$. The variation bound is spectral
Cauchy--Schwarz applied to the two quotient vectors of $A_1,A_2$.
The factors $\sqrt2$ from the diagonal coordinates have been absorbed
into $E',p'$.

For $h\in L^1(0,\infty)$ write
$\widetilde h(z)=\int_0^\infty h(a)e^{-az}\,\dd a$.
Integrating the two representations against
$h(a)\overline{h(c)}$ yields the decisive bound
\begin{equation}\label{eq:two-Laplace-bound}
 \int e^{-2DE}|\widetilde h(E-ip_i)|^2\,\dd\varrho
 \le\|A\|_2^2
 \sup_{u\ge0}|\widetilde h(u)|
 \sup_{v\in\R}|\widetilde h(iv)|.
\end{equation}
All integrations are justified by the finite variation of the
spectral measures and the $L^1$ assumption on $h$.

We choose transforms whose two axis suprema have bounded product,
but whose modulus grows exponentially outside $|\arg z|\le\pi/4$.
For $s=\pm1$ and $n\ge1$ use the principal-power functions
\begin{equation}\label{eq:cone-Laplace-tests}
 \widetilde h_n(z)
 =e^{-n}z^{1-4isn/\pi}(1+z)^{-3},\qquad
 \widetilde h_n(0)=0.
\end{equation}
Their suprema on the positive and imaginary axes are at most
$e^{-n}$ and $e^n$, respectively. Each is the Laplace transform of
an $L^1$ function. Here is the factorization that verifies this
point, including the boundary of the half-plane. For
$\nu=1-4isn/\pi$,
\[
 z^\nu(1+z)^{-3}
 =\bigl(1-(1+z)^{-1}\bigr)^\nu(1+z)^{-(3-\nu)}.
\]
The second factor is the transform of
$a^{2-\nu}e^{-a}/\Gamma(3-\nu)$, an integrable function.
The binomial coefficients of the first factor are absolutely
summable, since $|\binom\nu k|=O(k^{-2})$. It is therefore the
transform of the finite complex measure
\[
 \delta_0+
 \sum_{k\ge1}(-1)^k\binom\nu k
 \frac{a^{k-1}e^{-a}}{(k-1)!}\,\dd a.
\]
Convolution gives $h_n\in L^1$. Absolute convergence and
$\operatorname{Re}\nu=1$ also give the stated boundary value at
$z=0$.

Fix $\delta>0$ and the sector
$S_\delta=\{z\ne0:\operatorname{Re}z\ge0,
\ s\arg z\ge\pi/4+\delta\}$. On this sector,
\[
 |\widetilde h_n(z)|^2
 \ge e^{8\delta n/\pi}|z|^2|1+z|^{-6}.
\]
Substitution in \eqref{eq:two-Laplace-bound} gives
\[
 e^{8\delta n/\pi}
 \int_{S_\delta}e^{-2DE}|z|^2|1+z|^{-6}\,\dd\varrho
 \le\|A\|_2^2.
\]
The integrand is positive on $S_\delta$, so letting $n\to\infty$
shows $\varrho(S_\delta)=0$. Taking both signs and rational
$\delta>0$ gives $|\arg z|\le\pi/4$ on its support, or
$E\ge|p_i|$.

Finite linear combinations of characters with compact positive-side
tests lie in some such $B$ and have dense quotient classes in
$\mathcal H_j$. The spectral projection of $\{E<|p_i|\}$
vanishes on this dense set and hence vanishes identically.
\end{proof}

Applying this result to the nine-mirror laws gives
$H_j\ge|\mathsf P_i|$ for every ordered pair $i\ne j$. This coordinate
cone identifies the invariant projections used in clustering.

\section{Invariant observables and clustering}
\label{subsec:invariant-reflection}

Let $P_j$ be the orthogonal projection in $L^2(\lambda)$ onto the
observables fixed by all $U(te_j)$, and let $P$ project onto those
fixed by all translations. These are the corresponding conditional
expectations by \cref{lem:averages-translations}. The averages
$T^{-1}\int_0^TU(te_j)\,\dd t$ converge strongly to $P_j$.
The projections commute with one another and with every translation.

An observable fixed by translations along $e_j$ can be represented
in either half-space bounded by any plane $x_j=c$. Indeed, approximate
it in $L^2$ by finite sums of characters with compact tests, and translate
each approximant far enough into the chosen half-space. Translation
fixes the limiting observable and preserves the approximation error.
The half-space $L^2$ subspace is closed.

The linear geometric reflection $\mathcal R_j$ restricts to a positive
self-adjoint involution on this invariant subspace. Positivity follows
from its inclusion in the positive-side space and reflection
positivity. Its spectrum is therefore $\{1\}$, so
\begin{equation}\label{eq:invariant-reflection-fixed}
 \mathcal R_jA=A\qquad(A\in\operatorname{ran}P_j).
\end{equation}
Thus the quotient map is an isometry there. Its range is precisely
$\ker H_j$. To identify it, average
\[
 \langle[A]_j,e^{-tH_j}[B]_j\rangle
     =\langle\mathcal R_jA,U(te_j)B\rangle_{L^2(\lambda)}
\]
over $0\le t\le T$. Spectral dominated convergence on the left and
\cref{lem:averages-translations} on the right give
\begin{equation}\label{eq:zero-energy-class}
 P_0[B]_j=[P_jB]_j,\qquad P_0=\mathbf1_{\{0\}}(H_j).
\end{equation}
The image of $\operatorname{ran}P_j$ is closed because the map is
isometric, and contains the dense subspace $P_0[\mathcal E_j^+]_j$
of $\ker H_j$. This proves the asserted identification.

For a law invariant and reflection positive at the nine mirrors,
the cone forces the joint spectral measure at zero energy to be
concentrated at zero momentum. Hence $V_j(y)$ fixes $\ker H_j$.
If $A\in\operatorname{ran}P_j$ and $y\perp e_j$, then
\[
 [U(y)A-A]_j=0.
\]
The difference is again in $\operatorname{ran}P_j$, so the isometry
gives $U(y)A=A$. Consequently
\begin{equation}\label{eq:invariant-projections-equal}
 P_1=P_2=P_3=P.
\end{equation}

\subsection{Correlations in all directions}

\begin{proposition}\label{prop:all-direction-mixing}
Suppose $\lambda$ is translation invariant and invariant and
reflection positive at the nine mirrors. For all $A,B\in L^2(\lambda)$,
\begin{equation}\label{eq:all-direction-mixing}
 \langle A,U(a)B\rangle\longrightarrow\langle A,PB\rangle
                      \quad\text{as }|a|\longrightarrow\infty.
\end{equation}
If $\lambda$ is translation ergodic, the limit is
$\overline{\E_\lambda A}\,\E_\lambda B$ and
$\ker H_j=\C\Omega_j$ for every $j$.
\end{proposition}

\begin{proof}
First fix $j$. For $A=\mathcal R_jA'$ and $A',B\in\mathcal E_j^+$,
\eqref{eq:zero-energy-class} gives, for $t\ge0$ and $y\perp e_j$,
\begin{align*}
 |\langle A,U(te_j+y)(1-P_j)B\rangle|
 &=|\langle[A']_j,V_j(y)e^{-tH_j}(1-P_0)[B]_j\rangle|\\
 &\le\|[A']_j\|\,\|e^{-tH_j}(1-P_0)[B]_j\|.
\end{align*}
The last norm tends to zero by spectral dominated convergence;
the bound is uniform in $y$. Approximate general $A,B$ by local
character sums. Translate their supports into opposite half-spaces,
at a cost of replacing $t$ by $t-2r$ for one fixed $r$.
The correlation error is bounded uniformly in $t,y$ by
$\|A-A_1\|_2\|B\|_2+\|A_1\|_2\|B-B_1\|_2$.
The adjoint identity gives the same conclusion as $t\to-\infty$.
Thus
\[
 \lim_{|t|\to\infty}\sup_{y\perp e_j}
 |\langle A,U(te_j+y)(1-P_j)B\rangle|=0.
\]
By \eqref{eq:invariant-projections-equal}, $P_j=P$ for all $j$.
For any sequence $|a|\to\infty$, a subsequence has one fixed
coordinate tending in modulus to infinity. The displayed uniform
bound applies with that coordinate, while $U(a)PB=PB$.
Every subsequence has such a further subsequence, proving
\eqref{eq:all-direction-mixing}. In an ergodic law $P$ is the
projection onto constants. The isometry onto $\ker H_j$ then gives
its one-dimensionality.
\end{proof}

\section{Ergodic laws and inherited reflection forms}
\label{subsec:component-hypotheses}

We recover the reflection properties of conditional laws by
testing the parent reflection forms against invariant weights.

Let $\mathcal I$ be the translation-invariant event sigma algebra.
The evaluation space is standard Borel by \cref{lem:standard-borel}.
The regular conditional distribution theorem
\cite[Theorems~4.1.16--4.1.17]{Durrett2019} therefore gives conditional laws $\lambda_F$ given $\mathcal I$, with
\begin{equation}\label{eq:component-kernel}
 \lambda_F(\Phi)=\E_\lambda[\Phi\mid\mathcal I](F),\qquad
 \lambda=\int\lambda_F\,\dd\lambda(F).
\end{equation}
The first equality is understood almost surely for each integrable
$\Phi$ (and for nonnegative $\Phi$ by monotone convergence).
The image measure on laws is the translation-ergodic decomposition.
We recall why a single full-measure set of ergodic laws can be used.

For an open set $O$, $\Sspace(\overline O;V)$ denotes
the subspace of $\Sspace(\R^3;V)$ consisting of functions
supported in $\overline O$, with the induced Schwartz topology.
Fix a countable real test space $\mathcal D_0\subset C_c^\infty(\R^3;V)$
over $\mathbb Q$, stable under rational translations and cubic
isometries, with $\mathcal D_0\cap C_c^\infty(O;V)$ dense
in $\Sspace(\overline O;V)$ for $\R^3$ and the six coordinate
half-spaces $O$. Take the rational span of translated cubic images
of countable Schwartz-dense compact-test families in these seven
open sets to obtain it.
Let $\mathcal A_0$ consist of finite sums of $e^{iF(f)}$,
$f\in\mathcal D_0$, with coefficients in $\mathbb Q+i\mathbb Q$.
This same countable family is dense in $L^2$ for every law on the
evaluation sigma algebra and determines such laws. For density,
condition an observable orthogonal to its span on finitely many
evaluations. Rational combinations of the tests show that the
Fourier transform of the resulting finite complex measure vanishes
on rational frequencies. Continuity and Fourier uniqueness make it
zero. Conditional expectations on the increasing finite evaluation
sigma algebras converge to the observable, proving density.
The same finite-dimensional Fourier argument proves determination
of laws. For the elements
$A_p$, choose cube-average times $T_k\to\infty$ so that
\[
 \|M_{T_k}A_p-PA_p\|_2^2\le2^{-k}\quad(p\le k),\qquad
 M_T=T^{-3}\int_{[0,T]^3}U(a)\,\dd a.
\]
Equation \eqref{eq:component-kernel} and summability give convergence
in $L^2(\lambda_F)$ for every $p$ on one full-measure set of $F$.
Moreover $PA_p$ has zero conditional variance, and thus equals
the constant $\lambda_F(A_p)$ in that space. The conditional laws
are translation invariant: test the parent invariance against
$\mathcal I$-measurable bounded weights, first for rational
translations and the countable characters, then use continuity and
Fourier uniqueness. \Cref{lem:averages-translations} for each such law now
shows that its invariant projection is the constant projection,
first on the dense algebra and then on all $L^2(\lambda_F)$.
Thus these conditional laws are ergodic.

To prove uniqueness, the summable parent errors above give
$M_{T_k}A_p\to PA_p$ almost surely, simultaneously for all $p$.
Let $\lambda=\int\eta\,\dd\pi(\eta)$ be any decomposition into
translation-invariant ergodic laws. This parent full-measure event
has full $\eta$-measure for $\pi$-almost every $\eta$.
On the other hand, \cref{lem:averages-translations} gives
$M_{T_k}A_p\to\eta(A_p)$ in $L^2(\eta)$.
Uniqueness of the limit in probability therefore gives
$PA_p(F)=\eta(A_p)$ for $\eta$-almost every $F$, simultaneously
for all $p$. The conditional-law identity in
\eqref{eq:component-kernel} holds on another parent full-measure
event for this same countable family. Thus
$\lambda_F(A_p)=\eta(A_p)$ there for every $p$.
Since $\mathcal A_0$ determines laws, $\lambda_F=\eta$.
Taking the distribution of $\lambda_F$ under the mixture shows
that its image measure is $\pi$, proving uniqueness.

Fix an axis mirror. A bounded invariant weight $h$ satisfies
$\mathcal R_jh=h$ by \eqref{eq:invariant-reflection-fixed} and has a
representative in either half-space. If $h\ge0$ and
$A\in\mathcal E_j^+$, reflection positivity applied to $\sqrt hA$
gives
\begin{equation}\label{eq:weighted-component-rp}
 \E_\lambda[h(\Theta_jA)A]\ge0.
\end{equation}
Reflection invariance similarly gives
$\E_\lambda[h\mathcal R_jA]=\E_\lambda[hA]$.
Taking all bounded $\mathcal I$-measurable weights in these identities
proves that the corresponding conditional reflection form is
nonnegative and the conditional character expectations are reflection
invariant. Countably many positive-side character sums suffice;
density extends the statements to the full positive-side algebra
in each conditional law. Translation invariance moves the mirror
to every real height.

This argument also transfers any further fixed mirror with unit normal $n$.
Choose $Q\in O(3)$ with $Qe_1=n$ and apply the axis argument to
the transported law $(Q^{-1})_*\lambda$. The transported law is
reflection positive at $x_1=0$; it need not equal $\lambda$.
The map $Q_*$ sends translation-ergodic laws to translation-ergodic
laws and transports their reflection forms. By uniqueness of the
ergodic decomposition, its image decomposition is that of $\lambda$.
Applying this to the six diagonal mirrors, and intersecting the
resulting nine full-measure sets, transfers all required mirrors.

The nine reflections generate the cubic group, so these components
are cubically invariant. \Cref{prop:all-direction-mixing} applies
to each component and gives clustering and
$\ker H_j=\C\Omega_j$. Any specified measurable set of full parent
measure has full measure in almost every component by
\eqref{eq:component-kernel}. This proves all assertions of
\cref{thm:components-clustering} except the optional Gaussian bound.

\section{Gaussian bounds in the components}\label{sec:component-gaussian}
Suppose now that the parent law satisfies
\eqref{eq:probability-G} with the constant $C\ge0$ stated there. To transfer this Gaussian domination with
the same constant, fix real
$f\in\Sspace$ and put $G=e^{F(f)}$, $c=e^{C\|f\|_2^2/2}$,
$\phi=PG$. All moments of $G$ are finite. For each integer $n\ge2$,
Tonelli's theorem and the parent bound give
\begin{equation}\label{eq:component-exponential-average}
 \E_\lambda[(M_TG)^n]\le T^{-3n}
 \int_{([0,T]^3)^n}
 \exp\left(\frac {Cn}2\|f\|_2^2+
       \frac C2\sum_{p\ne q}\langle f,\tau_{a_p-a_q}f\rangle\right)\,\dd a.
\end{equation}
The overlaps tend to zero as $|a_p-a_q|\to\infty$ and are bounded
by $\|f\|_2^2$. Fix $\epsilon>0$. Outside the set where one pair is within distance
$\delta$, the right side is at most
$c^n e^{Cn(n-1)\epsilon/2}$ once $\delta$ is sufficiently large.
The exceptional set has relative volume at most
$\binom n2(2\delta)^3T^{-3}$ and its integrand is at most
$e^{Cn^2\|f\|_2^2/2}$.
Since $M_TG\to\phi$ in $L^2$, and the averages and $\phi$ are
uniformly bounded in $L^{2n}$, interpolation gives convergence in
$L^n$. Sending $T\to\infty$ and then $\epsilon\to0$ in
\eqref{eq:component-exponential-average} yields
$\|\phi\|_n\le c$ for every $n$, hence $\phi\le c$ almost surely.
By \eqref{eq:component-kernel},
\begin{equation}\label{eq:component-G}
 \lambda_F(e^{F(f)})\le e^{C\|f\|_2^2/2}\quad\text{almost surely}.
\end{equation}
Choose a countable Schwartz-dense real test family and intersect
these full-measure sets. For any other test, pointwise convergence
of the evaluations under each conditional law and Fatou's lemma
extend \eqref{eq:component-G} to it. Thus all real tests obey the
bound on one full-measure set of component laws.

This proves the Gaussian assertion of
\cref{thm:components-clustering}. For a two-form field, closedness
also passes to the components whenever it holds in the parent law:
the event $dF=0$ is measurable by testing it on a countable dense
family of three-form tests.
\appendix
\section{Measurability and translation averages}\label{sec:background}
We give the measure-theoretic details used in the component argument.
They require no moment bound.

\begin{lemma}\label{lem:standard-borel}
The space \(\mathcal X=\mathcal S'(\R^3;V)\), with its evaluation
sigma-algebra, is standard Borel. Translation is a jointly measurable
action on this space.
\end{lemma}
\begin{proof}
Choose a countable rational vector space \(\mathcal D_0\) dense in
\(\mathcal S(\R^3;V)\), and increasing norms \(q_n\) defining the
Schwartz topology. For \(n\ge1\), put
\[
 X_n=\{F:|F(f)|\le n q_n(f)\text{ for every }f\in\mathcal S\}.
\]
Continuity of each distribution gives \(\mathcal X=\bigcup_nX_n\).
The evaluation map
\(\iota(F)=(F(f))_{f\in\mathcal D_0}\) is injective into the
Polish space \(\R^{\mathcal D_0}\).
Its image on \(X_n\) is exactly the set of sequences satisfying
\[
 z_{sf+tg}=sz_f+tz_g,\qquad |z_f|\le n q_n(f)
 \quad(f,g\in\mathcal D_0,\ s,t\in\Q).
\]
These are closed conditions on countably many coordinates.
Conversely, the bound extends such a rational linear functional
continuously and uniquely to the Schwartz space; density and
continuity make the extension real linear.
Thus \(\iota(\mathcal X)\) is a countable union of closed sets.
Every evaluation is measurable with respect to the
\(\mathcal D_0\) coordinates by approximation and samplewise
continuity. Hence \(\iota\) is a measurable isomorphism onto a
Borel subset of a Polish space, which is standard Borel
\cite[Theorem~2.1.22]{Durrett2019}.

For fixed \(f\), the function
\((a,F)\mapsto F(\tau_{-a}\cdot f)\) is continuous in \(a\) and
measurable in \(F\). Replace \(a\) by its coordinatewise lower
dyadic approximation \(2^{-n}\lfloor2^na\rfloor\).
The resulting functions are measurable on the product space and
converge pointwise to the original one. Since evaluations generate
the sigma-algebra, the translation action is jointly measurable.
\end{proof}

\begin{lemma}\label{lem:averages-translations}
Let \(\lambda\) be translation invariant. If \(J\subset\{1,2,3\}\)
is nonempty, the averages
\[
 M_T^J A=T^{-|J|}\int_{[0,T]^J}U(a)A\,\dd a
\]
converge in \(L^2(\lambda)\) to the orthogonal projection \(P_J\)
onto the observables fixed by the translations in the coordinate
subspace spanned by \(\{e_j:j\in J\}\).
This projection is conditional expectation onto the corresponding
invariant sigma-algebra. All \(M_T^J\) are contractions on
\(L^p(\lambda)\), \(1\le p<\infty\).
\end{lemma}
\begin{proof}
Joint measurability and Fubini give the averages, and Jensen's
inequality gives the \(L^p\) contraction. The representation \(U\)
is strongly continuous by bounded convergence on local characters
and their density in \(L^2(\lambda)\).
Let \(\mathcal K\) be the closed span of
\(B-U(b)B\), where \(b\) is in the coordinate subspace.
Its orthogonal complement is the invariant subspace.
The averages fix that subspace. The symmetric-difference bound
for a cube and its translate gives
\[
 \|M_T^J(B-U(b)B)\|_2
 \le\frac{2|b|_1}{T}\|B\|_2.
\]
Thus the averages tend to zero on a dense subspace of
\(\mathcal K\), and their uniform norm bound extends this
convergence to all of \(\mathcal K\).
This proves the asserted orthogonal-projection limit.

The invariant subspace equals the \(L^2\) space of the invariant
sigma-algebra: indicators of invariant events are fixed, and
the inverse images of Borel sets under a fixed observable are
invariant modulo null sets. Therefore its orthogonal projection
is the stated conditional expectation. The coordinate projections
commute, since their averaging operators commute.
\end{proof}
\printbibliography
\end{document}
